%------------------------------------------------------------------------------%
% COERCIVE SESQUILINEAR FORMS                                                  %
%------------------------------------------------------------------------------%
% File  : Coercive_Sesquilinear_Forms.tex                                      %
% Author: Klaus Hoeflinger, Beethovenstrasse 11, D-73117 Wangen                %
% Email : khoeflinger@t-online.de                                              %
%------------------------------------------------------------------------------%

%------------------------------------------------------------------------------%
%                       DOCUMENT CLASS and INCLUDE FILES                       %
%------------------------------------------------------------------------------%
\documentclass[11pt,twoside]{article}
\usepackage{geometry}
\geometry{paperheight=241mm,
          paperwidth=152mm,
          top=22mm,
          bottom=17mm,
          left=20mm,
          right=20mm,
          textwidth=112mm,
          textheight=202mm,
          headheight=10mm,
          headsep=3mm,
          footskip=9mm,
          includehead,
          includefoot,
          marginparsep=0mm,
          marginparwidth=0mm}

\renewcommand{\baselinestretch}{0.945}\normalsize

\AtBeginDocument{\setlength\abovedisplayskip{2mm}}
\AtBeginDocument{\setlength\belowdisplayskip{2mm}}

%------------------------------------------------------------------------------%
% define title formats                                                         %
%------------------------------------------------------------------------------%
\usepackage{titlesec}

\renewcommand{\thesection}{\arabic{section}}
\renewcommand{\sfdefault}{FiraSans}
\renewcommand{\ttdefault}{pcr}
\renewcommand{\rmdefault}{antiqua}

\titleformat{\part}[display]
{\normalfont \large \rmfamily \bfseries \centering}
{\small{\thepart.}}{2pt}{}

\titleformat{\section}[runin]
{\normalfont \rmfamily \bfseries}
{\thesection}{1pt}{.\;}[.]

%\titleformat{\section}
%{\normalfont \bfseries \small \filcenter}
%{\thesection.\;}{0mm}{}

\titleformat{\subsection}[runin]
{\normalfont \itshape}
{}{0mm}{}

\titleformat{\subsubsection}[runin]
{\normalfont \itshape}
{}{}{}{}

\titleformat{\paragraph}[runin]
{\normalfont}
{}{}{}{}

\titleformat{\subparagraph}[runin]
{\normalfont}
{}{}{}{}

\titlespacing*{\part}          { 0pt}{ 0pt}{ 8pt}
\titlespacing*{\section}       { 0pt}{ 7pt}{ 4pt}
\titlespacing*{\subsection}    { 0pt}{14pt}{ 6pt}
\titlespacing*{\subsubsection} { 0pt}{ 6pt}{12pt}
\titlespacing*{\paragraph}     { 0pt}{ 6pt}{12pt}
\titlespacing*{\subparagraph}  { 0pt}{ 0pt}{12pt}

%------------------------------------------------------------------------------%
% define enumerate parameters                                                  %
%------------------------------------------------------------------------------%
\usepackage{enumitem}
\setlength{\parindent}{3pt}
\setlength{\parskip}  {0pt}

%\setlist{noitemsep}

\setlist[enumerate]{
         leftmargin=12pt,
         rightmargin=0pt,
         itemsep=1pt,
         itemindent=3pt,
         topsep=3pt,
         labelsep=4pt,
         partopsep=0pt,
         labelwidth=0pt,
         listparindent=0pt,
         parsep=0pt}

\setlist[itemize]{
         leftmargin=12pt,
         rightmargin=0pt,
         itemsep=0pt,
         itemindent=1pt,
         topsep=1pt,
         labelsep=4pt,
         partopsep=1pt,
         labelwidth=0pt,
         listparindent=0pt,
         parsep=0pt}

%------------------------------------------------------------------------------%
% define packages                                                              %
%------------------------------------------------------------------------------%
\usepackage{upref}
\usepackage{amssymb}
\usepackage{slantsc}
\usepackage{amsmath}
\usepackage{amsthm}
\usepackage{thmtools}
\usepackage{amscd}
\usepackage{verbatim}
\usepackage{latexsym}
\usepackage{multicol}
\usepackage{graphicx}
\usepackage{setspace}
\usepackage[ngerman,english]{babel}
\usepackage[protrusion=true,expansion=true]{microtype}
\usepackage{tikz}
\usetikzlibrary{arrows,arrows.meta,decorations.markings}
\usepackage{mathrsfs}
\usepackage{amsfonts}
\usepackage{datetime2}
\usepackage{textgreek}
\usepackage{hyperref}
\usepackage{pifont}
\usepackage{relsize}
%\usepackage{biblatex}
\relsize{-0.05}
\usepackage[skip=0mm]{parskip}
\usetikzlibrary{decorations.pathmorphing}

\usepackage{mlmodern}
\usepackage[T5,T1]{fontenc}
\usepackage{ragged2e}

%\usepackage{mathabx}
%\usepackage{times}

%\usepackage{fontspec}
%\setmainfont{Nimbus Roman No9 L}
%\usepackage[T1]{fontenc}

%------------------------------------------------------------------------------%
% define page layout                                                           %
%------------------------------------------------------------------------------%
\usepackage{fancyhdr}
\pagestyle{fancy}
\setlength{\headheight}{30.0pt}
\renewcommand{\headrulewidth}{0pt}

\AtBeginDocument{
  \addtolength\abovedisplayskip{-0.0\baselineskip}
  \addtolength\belowdisplayskip{-0.0\baselineskip}
}

\fancyhead{}
\fancyfoot{}
\addtolength{\headwidth}{0.02\marginparwidth}

\makeatletter
\let\ps@plain\ps@fancy
\makeatother

\renewcommand\sectionmark[1]
{
 \def\sectionname{#1}
 \markright{\thesection #1}
}

\makeatletter
\@addtoreset{section}{part}
\makeatother

%------------------------------------------------------------------------------%
% define footnote without number                                               %
%------------------------------------------------------------------------------%
\newcommand{\myfootnote}[1]{%
\renewcommand{\thefootnote}{}%
\footnotetext{#1}%
\renewcommand{\thefootnote}{\arabic{footnote}}%
}

%------------------------------------------------------------------------------%
% define numbering in equations                                                %
%------------------------------------------------------------------------------%
\numberwithin{equation}{section}

%------------------------------------------------------------------------------%
% define newtheorem                                                            %
%------------------------------------------------------------------------------%
\declaretheoremstyle[
headfont=\scshape, 
headindent = 0mm, %\parindent,
%postheadhook = {\hspace{0mm}\newline},
%postheadspace = \newline, 
spaceabove = 3mm, 
spacebelow = 3mm,
numberwithin=section,
name=Theorem]{khMATH}
\declaretheorem[style=khMATH]{theorem}

\declaretheoremstyle[
headfont=\bfseries, 
headindent = 0mm, %\parindent,
%postheadhook = {\hspace{0mm}\newline},
%postheadspace = \newline, 
spaceabove = 3mm, 
spacebelow = 3mm,
numberwithin=part,
name=Theorem]{khMATH1}
\declaretheorem[style=khMATH1]{theorem1}

\declaretheoremstyle[
headfont=\scshape, 
headindent = 0mm, %\parindent,
%postheadhook = {\hspace{0mm}\newline},
%postheadspace = \newline, 
spaceabove = 3mm, 
spacebelow = 3mm,
numberwithin=section,
name=Corollary]{khMATH1}
\declaretheorem[style=khMATH1]{corollary}

%            name         counter     label              numeration-mode
\theoremstyle{plain}
\newtheorem*{introduction}            {\sc Introduction}
\newtheorem {standard}                {}                 [section]
\newtheorem*{definition}              {\sc Definition}
\newtheorem{satz1}                    {\sc Satz}
\newtheorem{lemma}                    {\sc Lemma}
\newtheorem*{proposition}             {\sc Proposition}
%\newtheorem{corollary}               {\sc Corollary}
\newtheorem*{corollar}                {\sc Korollar}
\newtheorem*{remark}                  {\sc Remark}
\newtheorem*{remarks}                 {\sc Remarks}
\newtheorem*{example}                 {Example}
\newtheorem*{examples}                {Examples}
\newtheorem*{theo}                    {\sc Theorem}

%------------------------------------------------------------------------------%
% define tabular                                                               %
%------------------------------------------------------------------------------%
\usepackage{tabularx}
\newcolumntype{L}[1]{>{\raggedright\arraybackslash}p{#1}}
\newcolumntype{C}[1]{>{\centering\arraybackslash}  p{#1}} 
\newcolumntype{R}[1]{>{\raggedleft\arraybackslash} p{#1}} 

%------------------------------------------------------------------------------%
% define \sh, \zB                                                              %
%------------------------------------------------------------------------------%
\def\dsh{{d.\,h.}}
\def\ie{{i.\,e.}}
\def\zB{z.\,B.}
\renewcommand{\abstractname}{ }

%------------------------------------------------------------------------------%
% change default spacing for cases                                             %
%------------------------------------------------------------------------------%
\makeatletter
\renewcommand*\env@cases[1][1.6]{%
  \let\@ifnextchar\new@ifnextchar
  \left\lbrace
  \def\arraystretch{#1}%
  \array{@{}l@{\quad}l@{}}%
}
\makeatother

\newenvironment{rcases}
  {\left.\begin{aligned}}
  {\end{aligned}\right\rbrace}

%------------------------------------------------------------------------------%
% change default for \predisplaypenalty                                        %
%------------------------------------------------------------------------------%
\predisplaypenalty=990
\allowdisplaybreaks \numberwithin{equation}{section}

%------------------------------------------------------------------------------%
% general notations                                                            %
%------------------------------------------------------------------------------%
\def\*{\star}
\def\={\,=\,}
\def\+{\,+\,}
\def\xsubset{{\,\subset\,}}
\def\xtimes{{\,\times\,}}
\def\xeof{{\,\in\,}}
\def\eq{{\,=\,}}
\def\xto{{\,\to\,}}
\def\eqdef{\,:=\,}
\def\:{{\,:\,}}
\def\ele{{\,\in\,}}
\def\exists{\mbox{ exists }}
\def\and{\mbox{ and }}
\def\for{\mbox{ for }}
\def\fa{\mbox{ for all }}
\def\fe{\mbox{ for each }}
\def\qfa{\quad \mbox{ for all }}
\def\df{\,\dotfill}
\def\iff{if and only if }
\def\wrt{with respect to}

\makeatletter
\newcommand \Df {\leavevmode \cleaders \hb@xt@ .70em{\hss .\hss }\hfill \kern \z@}
\makeatother

\def\sql{\mathfrak{a}}               % sesquilinear form a
\def\DF{B}                           % Dirichlet form a
%------------------------------------------------------------------------------%
% number sets and special elements                                             %
%------------------------------------------------------------------------------%
\def\P{\mathbb{H}}                   % half plane of $\C$
\def\J{I}                            % interval I
\def\N{{\mathbb{N}}}                 % unsigned integers
\def\Q{{\mathbb{Q}}}                 % rational integers
\def\Z{{\mathbb{Z}}}                 % integers
\def\R{{\mathbb{R}}}                 % real numbers
\def\Rn{{\mathbb{R}}^n}              % real numbers in Rn
\def\Rd{{\mathbb{R}}^d}              % real numbers in Rn
\def\C{{\mathbb{C}}}                 % complex numbers
\def\F{{\mathbb{K}}}                 % arbitrary field
\def\Torus{{\mathbb{T}}}             % complex circle line
\def\T{\tau}                         % upper bound of interval (0,t)
\def\sign#1{\operatorname{sign}(#1)} % sign of a number

%------------------------------------------------------------------------------%
% set theory                                                                   %
%------------------------------------------------------------------------------%
\def\G{G}                            % domain $G$
\def\clG{\overline{\G}}              % closure of $G$
\def\bndG{\partial \G}               % boundary of $G$
\def\setA{\mathcal{A}}               % set A
\def\setB{\mathcal{B}}               % set B
\def\setC{\mathcal{C}}               % set C
\def\setM{M}                         % set M
\def\setN{N}                         % set N
\def\setS{\mathcal{S}}               % set S
\def\famF{\mathfrak{F}}              % family of sets
\def\indI{\mathcal{I}}               % index set I
\def\pot#1{\mathfrak{P}(#1)}         % power set

\def\relR{\mathbf{R}}                % relation R
\def\mapf{f}                         % mapping f
\def\mapg{g}                         % mapping g

\def\set#1#2{\{ \, #1 \,: \, #2 \, \}}
\def\tensor#1#2{#1 \otimes #2}
\def\Tens{\oplus}

\def\cal#1{{\mathcal{#1}}}
\def\aut#1{\operatorname{Aut}(#1)}

\def\cross#1#2{#1 {\,\times\,} #2}
\def\multicross#1#2{#1 \times  \ldots \times  #2}
\def\multitens#1#2 {#1 \otimes \ldots \otimes #2}

%------------------------------------------------------------------------------%
% group theory                                                                 %
%------------------------------------------------------------------------------%
\def\1{{\mathsf{1}}}                % unit element of a monoid
\def\0{{\mathsf{0}}}                % unit element of an Abelian monoid
\def\kernel#1{\mathfrak{N}(#1)}     % kernel of a homomorphism
\def\cokernel#1{\mathfrak{C}{(#1)}} % cokernel of a homomorphism
\def\Hom#1{\operatorname{Hom}(#1)}

%------------------------------------------------------------------------------%
% linear algebra                                                               %
%------------------------------------------------------------------------------%
\def\vecV{\mathit{V}}               % vector space V
\def\vecH{\mathit{H}}               % vector space H
\def\vecW{\mathit{W}}               % vector space W
\def\vecU{\mathit{U}}               % subspace U
\def\convC{\mathcal{C}}             % convex set C
\def\basH{\mathfrak{B}}             % Hamel base B

\def\LO#1{\mathscr{L}(#1) }         % algebra of linear transformations

\def\scalar#1#2{ \langle #1,#2 \rangle }
\def\trace#1{\operatorname{tr}(#1)}
\def\dim#1{{\operatorname{dim} \, #1 }}
\def\codim#1{{\operatorname{codim} \, #1 }}
\def\dim{\operatorname{dim}}
\def\codim{\operatorname{codim}}
\def\ind{\operatorname{ind}}

\def\genG{\cal{G}}

\def\algA{\cal{A}}
\def\algB{\cal{B}}
\def\algU{\cal{U}}

\def\idI{\cal{I}}

%------------------------------------------------------------------------------%
% topology                                                                     %
%------------------------------------------------------------------------------%
\def\compK{{K}}          % compact set C
\def\openO{\mathcal{O}}  % open set O
\def\openU{\mathcal{U}}  % open set U
\def\U{\mathcal{U}}      % neighbourhood U
\def\V{\mathcal{V}}      % neighbourhood V
\def\d{\mathit{d}}       % metric function d

\def\interior#1{\mathring{#1}}
\def\closure#1{{\overline{#1}}}

\def\topT{\mathcal{T}}   % topology T
\def\topX{X}             % topological space X
\def\topY{Y}             % topological space Y
\def\topZ{Z}             % topological space Z
\def\metrS{M}            % metric space M

\def\grpG{G}             % topological group G
\def\grpA{A}             % topological group A
\def\grpB{B}             % topological group B
\def\grpC{C}             % topological group C
\def\grpH{H}             % topological group H
\def\grpU{U}             % topological group U
\def\grpV{V}             % topological group V
\def\topG{G}             % topological group G
\def\topH{H}             % topological group H
\def\topU{U}             % topological subgroup U

\def\subS{\cal{S}}
\def\riR{R}              % ring R

%------------------------------------------------------------------------------%
% calculus                                                                     %
%------------------------------------------------------------------------------%
\def\supp#1{\operatorname{supp}(#1)}
\def\singsupp#1{\operatorname{sing \, supp}(#1)}

\def\re{\operatorname{Re}}
\def\im{\operatorname{Im}}
\def\grad{\operatorname{grad}}

%\def\abs#1 {\left|#1\right|}
\def\abs#1 {\lvert #1 \rvert}
%\def\abs#1 {\left|\,#1\,\right|}

%------------------------------------------------------------------------------%
% measure theory                                                               %
%------------------------------------------------------------------------------%
\def\salgA{\Sigma}               % sigma-algebra S
\def\salgB{\cal{B}}              % sigma-algebra B
\def\Borel#1{\cal{B}_{#1}}       % Borel-algebra 
\def\LB{\lambda}                 % Lebesgue-Borel measure

%------------------------------------------------------------------------------%
% function spaces                                                              %
%------------------------------------------------------------------------------%
\def\Lp{L^p(\G)}                 % spaces L_p
\def\Lq{L^q(\G)}                 % spaces L_q
\def\Lh{L^2(\G)}                 % spaces L_2
\def\Li{L^{\infty}(\G)}          % spaces L_infty
\def\Lc{L_{loc}^1(\G)}           % spaces L_1^{loc}
\def\Ck{C^k(\G)}                 % spaces C_k

\def\BV#1{BV(#1) }               % set of functions of bounded variation
\def\AC#1{AC(#1) }               % set of absolutely continuous functions

\def\MP#1{\mathscr{M}^+(#1) }    % set of positive measures
\def\MS#1{\mathscr{M}^-(#1) }    % vector space of finite signed measures
\def\MC#1{\mathscr{M}(#1) }      % vector space of complex measures
\def\MCR#1{\mathscr{M}_{r}(#1) } % vector space of regular complex measures

\def\SobW{W^{k,p}(\G)}                % Sobolev space W
\def\SobO{\mathring{W}^{k,p}(\G)}     % Sobolev space W
%\def\WN#1#2{\mathring{W}^{#1}(#2) }   % Sobolev space W
\def\WN#1#2{W_0^{#1}(#2) }   % Sobolev space W
%\def\HN#1#2{\mathring{H}^{#1}(#2) }   % Sobolev space H
\def\HN#1#2{H_0^{#1}(#2) }   % Sobolev space H
%------------------------------------------------------------------------------%
% functional analysis                                                          %
%------------------------------------------------------------------------------%
\def\snorm#1{\lVert #1 \rVert}
\newcommand{\norm}[1]{\left\lVert #1 \right\rVert}

\def\topV{V}                       % topological vector space V
\def\topW{W}                       % topological vector space W
\def\normS{X}                      % normed space

\def\banX{\mathit{X}}              % Banach space X
\def\banY{\mathit{Y}}              % Banach space Y
\def\banZ{\mathit{Z}}              % Banach space Z
\def\dbanX{\mathit{X^*}}           % dual space of Banach space X
\def\dbanY{\mathit{Y^*}}           % dual space of Banach space Y
\def\dbanZ{\mathit{Z^*}}           % dual space of Banach space Z

\def\banA{\mathit{A}}              % Banach algebra A
\def\banB{\mathit{B}}              % Banach algebra B
\def\banC{\mathit{C}}              % C*-algebra C


\def\domain#1{\mathfrak{D}(#1)}    % domain
\def\range#1{\mathfrak{R}(#1)}     % range

\def\isoJ{\mathfrak{J}}            % isometry J
\def\repA{{\cal{A}_{\sql}}}        % representation operator A
\def\linS{{S}}                     % linear operator S
\def\linG{{G}}                     % linear operator G
\def\linL{{L}}                     % linear operator L
\def\linT{{T}}                     % linear operator A
\def\linA{\mathit{A}}              % linear operator A
\def\linB{{B}}                     % linear operator B
\def\linM{{M}}                     % linear operator M
\def\linU{{U}}                     % linear operator U
\def\linI{{I}}                     % linear operator I
\def\A{\cal{A}}                    % linear differential operator A
\def\BDO{\cal{B}}                  % linear boundary differential operator B
\def\linU{{U}}                     % unitary operator U
\def\linP{{P}}                     % projection P
\def\compA{k}                      % compact operator k
\def\linFR{f}                      % finite rank operator f
\def\fredA{A}                      % Fredholm operator F

\def\HAM{\cal{H}}                  % Hamilton operator
\def\PA{\partial^{\alpha}}         % elementary differential operator
\def\DO{\cal{A}}                   % linear differential operator
\def\BC{\cal{B}}                   % boundary operator
\def\SLO{\cal{L}_{p,q}}            % Sturm-Liouville operator
\def\MO{\cal{L}_{M}}               % momentum operator

\def\HS#1{S(#1)           }        % algebra of Hilbert-Schmidt operators
\def\FR#1{F(#1)           }        % algebra of finite rank operators
\def\TC#1{N(#1)           }        % algebra of trace class operators
\def\BU#1{\mathscr{U}(#1) }        % algebra of unitary linear operators
\def\BO#1{\mathscr{L}(#1) }        % algebra of bounded linear operators
\def\CO#1{\mathscr{C}(#1) }        % algebra of compact linear operators
\def\CL#1{\mathscr{C}(#1) }        % algebra of compact linear operators
\def\FO#1{\Phi(#1) }               % set of Fredholm operators
\def\FK#1#2{\Phi_{#2}(#1) }        % set of Fredholm operators with index k

\def\hilH{\mathit{H}}              % Hilbert space H
\def\clU{\mathit{U}}               % closed subspace U

\def\orthO{\mathfrak{O}}           % orthonormal system O
\def\orthS{\mathfrak{S}}           % orthonormal system S
\def\basB{\mathfrak{B}}            % basis B of a Hilbert space

\def\GL#1 {\mathbf{GL}(#1)}        % linear group
\def\OG#1 {\mathbf{O}(#1)}         % linear group
\def\SOG#1 {\mathbf{SO}(#1)}       % linear group

%------------------------------------------------------------------------------%
% distribution theory                                                          %
%------------------------------------------------------------------------------%
\def\disF{F}                  % distribution F
\def\disG{G}                  % distribution G
\def\disH{H}                  % distribution H
\def\disU{U}                  % distribution U
\def\disT{T}                  % distribution T
\def\SF{C^{\infty}(\G)}       % space of smooth functions
\def\TF{C_c^{\infty}(\G)}     % space of compactly supported smooth functions
\def\TFRd{C_c^{\infty}(\R^d)} % space of test functions
\def\SRd{\mathscr{S}(\R^d)}   % Schwartz space
\def\SR1{\mathscr{S}(\R)}     % Schwartz space
\def\B1#1{C^1(#1) }           % space of continuously differentiable functions
\def\Bm#1{C^m(#1) }           % space of continuously differentiable functions
\def\Ck{C^k(\G)}
\def\TD{\mathscr{S}'(\R^d)}   % space of temperated distributions
\def\E{\cal{E}}               % space of distributions with compact support
\def\D{\mathscr{D}}           % D
\def\FT{\mathcal{F}}          % Fourier transformation F
\def\FT{\mathscr{F}}          % Fourier transformation F

%------------------------------------------------------------------------------%
% other definitions                                                            %
%------------------------------------------------------------------------------%
\def\Proof{\textsc{Proof}}
\def\FRECHET{Fr\'{e}chet}
\def\ARZELA{Arzel\'{a}}
\def\GARDING{G\r{a}rding}
\def\HOELDER{H\"older}
\def\SCHROEDINGER{Schr\"odinger}
\def\JOERGENS{J\"orgens}
\def\WUEST{W\"uest}
\def\KAEHLER{K\"aehler}
\def\HOEFLINGER{H\"oflinger}
\def\POINCARE{Poincar\'{e}}
\def\FA{Functional Analysis}

\def\Author   {K.\,{\HOEFLINGER}}
\def\AuthorUC {K.\,HOEFLINGER}
\def\Version  {1.0}
\def\Email    {khoeflinger@t-online.de}


%\renewcommand\thesection{\Roman{section}}

\def\Title {Coercive Sesquilinear Forms}

\titleformat{\section}
{\normalfont \bfseries \filcenter}
{\thesection.\;}{0mm}{}

\titleformat{\subsection}[runin]
{\normalfont \filcenter}
{}{0mm}{}

\titleformat{\subsubsection}[runin]
{\itshape \filcenter}
{}{0mm}{}

\titlespacing*{\section}      {5pt}{12pt}{ 4pt}
\titlespacing*{\subsection}   {0pt}{ 6pt}{12pt}
\titlespacing*{\subsubsection}{0pt}{ 6pt}{ 3pt}

%------------------------------------------------------------------------------%
%                                BEGIN OF DOCUMENT                             %
%------------------------------------------------------------------------------%
\begin{document}
\thispagestyle{empty}

\centerline{\large{\textbf{\Title}}}
\vspace{4pt}
\centerline{\small{{\textsc{\Author}}}}
\vspace{10pt}

\fancyhead[C] {\small{\textsc{\Title}}}
\fancyhead[OL]{\small{\thepage}}
\fancyhead[ER]{\small{\thepage}}

\myfootnote
{
  Version {\Version} from \today;
  Email:  \textit{khoeflinger@\,t-online.de}
}

%------------------------------------------------------------------------------%
%                                 INTRODUCTION                                 %
%------------------------------------------------------------------------------%
Our aim in this text
is to present the general theory of coercive sesquilinear forms
and its applications to stationary linear equations.
Sesquilinear forms are mostly efficient
when studying abstract problems in the Hilbert space setting.

%------------------------------------------------------------------------------%
\section{Generalities}                                                         %
%------------------------------------------------------------------------------%
By a \textit{sesquilinear form} defined in complex linear space $\vecV$
is meant a mapping $\sql\:\vecV \xtimes \vecV \xto \C$,
which is linear in the first and conjugate linear in the second argument,
that is,
for all $x,y,z \in \vecV$ and $\lambda,\mu \in \C$ the computation rules
\[
\begin{cases}[1.2]
\quad \sql(x+y,z) \= \sql(x,z) \, + \, \sql(y,z) \\
\quad \sql(x,y+z) \= \sql(x,y) \, + \, \sql(x,z) \\
\quad \sql(\lambda x,\mu y) \= \lambda \overline{\mu} \, \sql(x,y) \\
\end{cases}
\]
are valid.
If the linear space $\vecV$ is real, 
however,
then $\sql$ is assumed to be linear also in the second argument,
so it becomes into a \textit{bilinear form} in this case.

\subparagraph*{} %- quadratic form
The mapping $x \in \vecV \mapsto \sql[x] := \sql(x,x) \in \C$ is called
the \textit{quadratic form} of $\sql$.
Notice that all values $\sql(x,y)$ can be recovered
from the values $\sql[x]$ and $\sql[y]$ by certain polarization formulas.

\paragraph*{} %- representation operators
A sesquilinear form $\sql$ may also be regarded
as a linear mapping $\repA\:\vecV \xto \vecV_a^{\#}$,
where $\vecV_a^{\#}$ is the (algebraic) conjugate space of $\vecV$,
that is,
the linear space of conjugate linear functionals $\phi\:\vecV \xto \C$.
More precisely,
given a sesquilinear form $\sql$,
each element $x \in \vecV$ gives rise to a conjugate linear functional
$\phi_x\:\vecV \xto \C$,
defined by $\phi_x(y) := \sql(x,y)$ for $y \in \vecV$.
Let $\vecV_a^{\#}$ again be the linear space
of conjugate linear functionals on $\vecV$.
Then the mapping
\[
\repA: \;
\begin{cases}[1.2]
\quad \repA\:V \xto V^{\#},\\
\quad \repA \, x := \phi_x \quad (x \in \vecV)
\end{cases}
\]
proves to be linear and is said to be the \textit{form operator}
or the \textit{representation operator} associated with $\sql$.

\subparagraph*{}
Reversely,
for any linear operator $\cal{A}\:\vecV \xto \vecV_a^{\#}$,
there is a sesqui\-linear form $\sql\:\vecV \xtimes \vecV \xto \F$
such that $\sql(x,y) \eq (\cal{A} \, x) y$ for $x,y \in \vecV$.
Hence sesquilinear forms acting on a complex linear space $\vecV$
appear in a one-to-one correspondence to linear mappings
acting between the spaces $\vecV$ and $\vecV_a^{\#}$.
This correspondence remains valid,
if $\vecV$ is a \textit{real} linear space,
$\sql$ is a \textit{bilinear} form
and $\cal{A}$ is the associated mapping
from $\vecV$ to the adjoint space $\vecV^{\#}$.

%------------------------------------------------------------------------------%
\section{Continuous and Coercive Forms}                                        %
%------------------------------------------------------------------------------%
From now on we assume that $\vecV$ is a \textit{normed} linear space
with underlying field $\C$.
Then any sesquilinear form $\sql\:\vecV \xtimes \vecV \xto \C$
is called \textit{bounded},
if there is $\Theta > 0$ such that
\[
\abs{\sql(x,y)} \, \leq \Theta \norm{x} \cdot \norm{y}
\quad \quad (x,y \in \vecV).
\]
Due to the \textit{uniform boundedness principle},
this property is equivalent to say
that $\sql$ is continuous on the space $\vecV \xtimes \vecV$.

\subparagraph*{} %- representation of continuous forms
The prototypical example of a continuous form is
the inner product in a complex Hilbert space $\hilH$,
where boundedness is ensured by the \textsc{Cauchy-Schwarz} inequality
$\abs{ \scalar{x}{y} } \leq \norm{x} \cdot \norm{y}$.

\subparagraph*{}
Moreover,
for each continuous sesquilinear form $\sql$ on $\hilH$
there is a bounded linear operator $\linB_{\sql}\:\hilH \xto \hilH$
such that
\begin{equation}\label{FORM_REPRESENTATION}
\sql(x,y) \= \scalar{\linB_{\sql} x}{y} \qfa x,y \in \hilH,
\end{equation}
which is a representation for the form $\sql$.
In fact,
if $x \in \hilH$ is fixed,
then the mapping $\sql(x,\cdot)\:\hilH \xto \C$
causes a continuous antilinear functional $\phi_x \in \hilH_a^*$,
and by \textsc{Riesz}'s representation theorem
there is $\tilde{x} \in \hilH$ such that
$\phi_x(y) \eq \scalar{\tilde{x}}{y}$ for all $y \in \hilH$.
This suggests to define $\linB_{\sql} x := \tilde{x}$ for every $x \in \hilH$.
If \eqref{FORM_REPRESENTATION} holds for some operator $\linB \in \BO{\hilH}$,
then the form $\sql$ is said to be represented by $\linB$.

\subparagraph*{} %- coercive sesquilinear forms
Finally,
any continuous sesquilinear form $\sql\:\hilH \xtimes \hilH \xto \C$
is called \textit{coercive}
if there is $\theta > 0$ such that
\begin{equation}\label{COERCIVITY_PROPERTY}
\re \, \sql(x,x) \, \geq \theta \norm{x}^2
\qfa x \in \hilH.
\end{equation}
The estimate \eqref{COERCIVITY_PROPERTY} especially implies the validness of
\[
\norm{\linB_{\sql} x}\norm{x}    \, \geq \,
\abs{\scalar{\linB_{\sql} x}{x}} \, \eq  \,
\abs{\sql(x,x)}                  \, \geq \,
\re \sql(x,x)                    \, \geq \,
\theta \norm{x}^2
\]
for all $x \in \hilH$.
Hence the representation operator $\linB_{\sql}$ is bounded below
with lower bound $\theta$,
and we may conclude by the fundamental \textit{bounded-below theorem}
that it is injective and the range $\range{\linB_{\sql}}$
forms a closed linear subspace in $\hilH$.

%------------------------------------------------------------------------------%
\section{The Variational Approach}                                             %
%------------------------------------------------------------------------------%
Coercivity of a sesquilinear form $\sql$ is of great interest {\wrt}
the solvability of the so-called \textit{variational problem},
which may be formulated as follows:
given a continuous and conjugate linear functional $\phi \in \vecV_a^*$,
it is to find $x \in \vecV$ such that
\begin{equation}\label{VARIATIONAL_EQUATION}
\sql(x,y) \= \phi(y) \qfa y \in \vecV.
\end{equation}

\subparagraph*{} %- Lax-Milgram theorem
The \textsc{Lax-Milgram} \textit{theorem} is of outstanding meaning
in applied functional analysis
and ensures well-posedness of the variational problem
for any right-hand side $\phi \in \vecV_a^*$:

\begin{theorem}\label{LAX-MILGRAM_THEOREM}
If $\sql\:\vecV \xtimes \vecV \xto \C$ is bounded and coercive,
then for each $\phi \in \vecV_a^*$
the equation \eqref{VARIATIONAL_EQUATION} owns exactly
one solution $x_{\phi} \in \vecV$.
In addition,
$\norm{\,x_{\phi}\,} \leq \theta^{-1} \norm{\phi}$ holds.
\end{theorem}

\subparagraph*{}
In other words,
if $\sql$ is bounded and coercive,
then the form operator $\repA\: \vecV \xto \vecV_a^*$ is bjective.
The bounded inverse theorem implies
that the inverse operator $\repA^{-1}$ is bounded,
too,
so the solution $x_{\phi}$ of \eqref{VARIATIONAL_EQUATION}
depends continuously on $\phi$.

\begin{comment}
\subparagraph*{}
The proof of theorem \ref{LAX-MILGRAM_THEOREM} is based on the fact
that the representation operator $\linB_{\sql}$ of $\sql$ is bounded below.

\Proof:
Since $\linB_{\sql}$ is bounded below,
it is injective and $\range{\linB_{\sql}}$ is closed.
It is still to show that $\linB_{\sql}$ is surjective.
If $z \in \range{\linB_{\sql}}^{\bot}$,
then for $\linB_{\sql} z \in \range{\linB_{\sql}} $ we have
\[
0 \=
\scalar{\linB_{\sql} z}{z}
\=
\sql(z,z)
\, \geq \,
\theta \norm{z}^2
\, \geq \, 0,
\]
which implies $z \eq 0$ and therefore $\overline{\range{\linB_{\sql}} } \eq \vecV$.
But $\range{\linB_{\sql}}$ is closed,
so $\range{\linB_{\sql}} \eq \vecV$.
Finally,
the inequality $\norm{x_f} \leq \theta^{-1} \norm{f}$
follows from the fact that $\norm{\linB_{\sql}^{-1}} \leq \theta^{-1}$.
\qed
\end{comment}

\paragraph*{} %- operator version of the \textsc{Lax-Milgram} theorem
There is also an \textit{operator version} of theorem \ref{LAX-MILGRAM_THEOREM}.
Let $\linA\:\vecV \xto \vecV$ be linear and bounded.
We shall denote $\linA$ \textit{coercive} if
\[
\re \scalar{\linA x}{x} \, > \, \theta \norm{x}^2
\]
holds for some $\theta > 0$ and all $x \in \vecV$.
It should be clear
that the mapping $(x,y) \mapsto \scalar{\linA x}{y}$
defines a coercive form on $\vecV$.
Then the \textsc{Lax-Milgram} theorem implies
invertibility of $\linA$ and boundedness of $\linA^{-1}$,
where $\norm{\linA^{-1}} \leq \theta^{-1}$ holds.

%------------------------------------------------------------------------------%
\section{Operators Associated with H-Elliptic Forms}                           %
%------------------------------------------------------------------------------%
It is often indicated to study abstract linear problems
in the setting of so-called \textit{pairs} $(\vecV,\hilH)$ of Hilbert spaces,
where the space $\vecV$ is continuously and densely embedded
into a second space $\hilH$.
More precisely,
we assume that

\begin{itemize}[leftmargin=9mm]
\item[(T-1)]
there is a linear mapping $\iota_1\:\vecV \xto \hilH$
such that $\overline{\range{\iota_1} }  \eq \hilH$;

\item[(T-2)]
there is $c>0$ such that
$\norm{\iota_1 v}_{\hilH} \leq c \norm{v}_{\vecV}$ for all $v \in \vecV$.
\end{itemize}

By means of the embedding $\iota_1$,
one can also think of $\vecV$ as a dense linear subspace in the space $\hilH$.

\subparagraph*{} %- Gelfand triples
Each Hilbert pair $(\vecV,\hilH)$ with embedding $\iota\:\vecV \xto \hilH$
gives rise to introduce the \textit{transposed pair} $(\hilH_a^*,\vecV_a^*)$,
where the conjugate space $\hilH_a^*$
can be seen as continuously and densely embedded
into the conjugate space $\vecV_a^*$
by a further embedding $\iota_2\:\hilH_a^* \xto \vecV_a^*$.
Then,
by identifying the spaces $\hilH$ and $\hilH_a^*$ via the \textsc{Riesz} map,
combining both Hilbert pairs leads to the triple
\begin{equation}\label{GELFAND_TRIPLE}
\vecV \, \overset{\iota_1}{\hookrightarrow} \,
\hilH \, \overset{\iota_2}{\hookrightarrow} \, \vecV_a^*,
\end{equation}
sometimes called a \textit{Gelfand triple},
an \textit{evolution triple} or a \textit{rigged Hilbert space}.
The center space $\hilH$ is often denoted
the \textit{pivot space} of \eqref{GELFAND_TRIPLE}.

\subparagraph*{}
Gelfand triples own a remarkable property.
Let $(\,\cdot\,,\,\cdot\,)$ and $\scalar{\,\cdot\,}{\,\cdot\,}$
be the inner products in $\vecV$ and $\hilH$,
respectively,
then for fixed chosen $x \in \hilH$
the action of the antilinear linear functional
$f_x := \iota_2 \, x \in \vecV_a^*$ to any element $v \in \vecV$
is represented by the inner product of $x$ and $v$ in $\hilH$,
{\ie} $f_x(v) \eq \scalar{x}{\iota_1 \, v}$ holds for all $v \in \vecV$.

\subsubsection*{} %- form operators
%----------------------------------
The presented theory of coercive sesquilinear forms in a Hilbert space $\vecV$
is often not sufficent for applications.
Hence it is quite natural to examine sesquilinear forms
also {\wrt} pairs $(\vecV,\hilH)$ of Hilbert spaces.
In applications it is often the case
that sesquilinear forms are merely continuous but not coercive.
In the setting of a Hilbert pair $(\vecV,\hilH)$,
however,
we try to make a form $\sql$ coercive
by adding an additional term,
which is just a multiplicity of the inner product in $\hilH$.

\subparagraph*{}
Indeed,
if $\sql\:\vecV \xtimes \vecV \xto \C$ is continuous
and the slightly modified form
\[
\sql_{\lambda}(u,v) \eqdef
\sql(u,v) \, + \, \lambda \scalar{\iota_1 u}{\iota_1 v}\, \quad
(u,v \in \vecV)
\]
is coercive for some $\lambda \in \R$,
that is,
there is $\theta > 0$ such that
\begin{equation}\label{QUASI_COERCIVE}
\re \sql(u,u)  \, + \,
\lambda \norm{\iota_1u}_{\hilH}^2 \, \geq \, \theta \norm{u}_{\vecV}^2
\quad (u \in \vecV),
\end{equation}
then $\sql$ is called
\textit{quasi-coercive},
\textit{$\hilH$-elliptic}
or simply a \textit{\GARDING} \textit{form}.
It should be clear that
if $\sql_{\lambda_0}$ is coercive for some $\lambda_0 \in \R$,
then $\sql_{\lambda}$ is so for $\lambda > \lambda_0$.

\paragraph*{} %- the operator associated with a quasi-coercive form
Quasi-coercive sesquilinear forms are often used
to define densely defined and closed linear operators
acting in the pivot space $\hilH$ of a given Gelfand triple.

\subparagraph*{}
Let $\sql$ be a sesquilinear form defined on a Hilbert space $\vecV$,
which is the left hand side of a Gelfand triple $(\vecV,\hilH,\vecV_a^*)$
with continuous embeddings
$\iota_1\:\vecV \xto \hilH$ and
$\iota_2\:\hilH \xto \vecV_a^*$.
It is our aim to introduce a linear operator $\linA_{\sql}$
with domain space $\hilH$ and range space $\hilH$.
This will be done by restricting the form operator
$\repA\:\vecV \xto \vecV_a^*$ associated with $\sql$
to those elements $u \in \vecV$,
for which $\repA \, u \in \hilH$.
Indeed,
the domain $\domain{\linA_{\sql}}$ is defined to be the subspace
of those elements $u \in \vecV$,
for which $\repA \, u \in \iota_2(\hilH) \subset \vecV_a^*$ holds,
briefly denoted
\[
\domain{\linA_{\sql}} \eqdef \{ \, u \in \vecV: \repA \, u \in \hilH \, \}.
\]

\subparagraph*{}
Clearly,
the set $\domain{\linA_{\sql}}$ can be regarded as a linear subspace
of $\hilH$ by means of the embedding $\iota_1$.
Then for $u \in \domain{A}$ there exists $y \in \hilH$
such that $\iota_2 y \eq \repA \, u$.
We set $x := \iota_1 u$ and $\linA_{\sql} x := y$
for $x \in \domain{\linA_{\sql}}$
and denote $\linA_{\sql}$
the \textit{operator in $\hilH$ associated with $(\vecV,\sql)$},
or shortly,
the \textit{variational operator} derived from $\sql$.
Another phrase is to say that
the operator $\linA_{\sql}$ is the \textit{part of $\repA$ in $\hilH$}.

\paragraph*{} %- main theorem on quasi-coercive forms
The following theorem collects a couple of nice properties 
of the operator $\linA_{\sql}$
associated with a continuous and quasi-coercive form $\sql$.
Besides the \textsc{Lax-Milgram} theorem,
is is our second important theorem in this text:

\begin{theorem}\label{QUASI-COERCIVE_FORMS_AND_ASSOCIATED_OPERATORS}
Let $(\vecV,\hilH,\vecV_a^*)$ be a triple of Hilbert spaces
and the sesquilinear form $\sql$ be continuous and quasi-coercive
on $\vecV$ with parameter $\lambda_0 \in \R$.
Then the following assertions are true:

\begin{itemize}[leftmargin=9mm]
\item[(a)]
$\linA_{\sql}$ is closed.

\item[(b)]
$\linA_{\sql}$ is densely defined in $\hilH$.

\item[(c)]
If the embedding of $\vecV$ into $\hilH$ is \textit{compact},
then the operator $\linA_{\sql} + \lambda$ has a compact inverse
$(\linA_{\sql} +\lambda)^{-1}$ for each $\lambda > \lambda_0$.

\item[(d)]
The operator $\linA_{\sql} + \lambda$ is m-accretive.

\item[(e)]
The operator $\linA_{\sql} + \lambda$ is sectorial.
\end{itemize}
\end{theorem}

\begin{comment}
\subparagraph*{}
To prove that $\linA_{\sql}$ is closed,
let $\{x_n\}_{n \in \N}$ be a sequence in $\domain{\linA_{\sql}}$
with limit element $x \in \hilH$
and $(\linA_{\sql} x_n$ be the sequence mapped by $\linA_{\sql}$
with limit element $y \in \hilH$.
It is to show that $x \in \domain{\linA_{\sql}}$ and $\linA_{\sql} x=y$.
Let $\abs{x} $ and $\norm{x}$ be the norms in $\vecV$ and $\hilH$,
respectively.
Clearly,
we have
\[
\norm{(\linA_{\sql}+\lambda) x_n - y - \lambda x} \xto 0 
\quad \for n \xto \infty.
\]
In the following we shall use the fact
that for $\lambda > \lambda_0$ the operators
$\linA_{\sql}+\lambda$ and $(\linA_{\sql}+\lambda)^{-1}$
forms linear homeomorphism between the normed spaces
$(\hilH,\norm{\,\cdot\,})$
and $(\domain{\linA_{\sql}},\abs{\,\cdot\,} )$,
so the latter one forms a closed subspace
within the normed space $(\vecV,\abs{\,\cdot\,} )$.
By this reason,
\[
x_n = (\linA_{\sql}+\lambda)^{-1} (\linA_{\sql}+\lambda) x_n
\to 
(\linA_{\sql}+\lambda)^{-1} (y + \lambda x)
\quad \for n \xto \infty
\]
{\wrt} the norm $\abs{\,\cdot\,} $ in $\domain{\linA_{\sql}}$.
Since $x$ belongs to $\hilH$ and $\domain{\linA_{\sql}}$ is closed,
the element $(\linA_{\sql}+\lambda)^{-1} x$
belongs to $\domain{\linA_{\sql}}$ with $\linA_{\sql} x \eq y$.

\subparagraph*{}
To prove that $\domain{\linA_{\sql}}$ is dense in $\hilH$,
suppose that $y \in \hilH$ fulfills $\scalar{x}{y} \eq 0$
for all $x \in \domain{\linA_{\sql}}$.
It is to show that $y=0$ in this case.
Since $\linA_{\sql}+\lambda$ is surjective,
there must exist $x_0 \in \domain{\linA_{\sql}}$ such that
$(\linA_{\sql}+\lambda) x_0 \eq y$.
Then
\[
0 \=
\scalar{\linA_{\sql} x_0}{x} \, + \, \scalar{\lambda x_0}{x} \=
\sql(x_0,x) \, + \, \lambda \scalar{x_0}{x}.
\]
Especially for $x \eq x_0$ the last equation leads to
\[
0 \=
\scalar{\linA_{\sql} x_0}{x_0} \, + \, \scalar{\lambda x_0}{x_0} \,\geq\,
\theta \norm{x_0}^2
\]
for some $\theta > 0$,
due to to quasi-coercivity of $\sql$,
hence $x_0 \eq 0$ and finally $y=0$.

\subparagraph*{}
Finally,
if $\iota_1 \: \vecV \xto \hilH$ is compact,
then $\iota_1 \circ (\linA_{\sql}+ \lambda)^{-1} \: \hilH \xto \hilH$
is compact,
too,
since the composition of a compact linear operator
and a continuous linear operator proves to be compact.
\qed.
\end{comment}

\paragraph*{} %- operators associated with adjoint and symmetric forms
It is often helpful to consider the \textit{adjoint form} $\sql^*$
of a given sesquilinear form $\sql$ defined by
\[
\sql^* (u,v) \eqdef \overline{\sql(v,u)}
\quad \quad (u,v \in \vecV).
\]
Notice that if $\sql$ is continuous,
then $\sql^*$ is so.
Moreover,
the form $\sql$ is called \textit{symmetric},
if $\sql \eq \sql^*$ holds\footnote{
\,A well-known example of a symmetric form is the inner product
$\scalar{\,\cdot\,}{\,\cdot\,}_{\hilH}$ in a Hilbert space $\hilH$,
which differs from arbitrary symmetric sesquilinear forms
only by positive definiteness.}.

\subparagraph*{} %- symmetric forms imply self-adjoint variational operators
With respect to the variational operators $\linA_{\sql}$ and $\linA_{\sql^*}$,
the following theorem is valid:

\begin{theorem}\label{QUASI-COERCIVE_FORMS_AND_ASSOCIATED_OPERATORS_II}
If $\sql^*$ is the adjoint form of $\sql$,
then ${\linA_{\sql}}^* \eq \linA_{\sql^*}$.
Furthermore,
if $\sql$ is symmetric,
then $\linA_{\sql}$ is self-adjoint.
\end{theorem}

\paragraph*{} %- generating operators
Given a Hilbert pair $(\vecV,\hilH)$,
the theorem \ref{QUASI-COERCIVE_FORMS_AND_ASSOCIATED_OPERATORS_II}
is very helpful to construct self-adjoint linear operators in $\hilH$
by starting with symmetric forms in $\vecV$.
This approach can especially be performed,
if the form is the inner product in $\vecV$.
The operator $\linA$ associated with the inner product
$\scalar{\,\cdot\,}{\,\cdot\,}_{\,\vecV}$,
called the \textit{generating operator} of $(\vecV,\hilH)$,
is coercive and even self-adjoint in $\hilH$
by theorem \ref{QUASI-COERCIVE_FORMS_AND_ASSOCIATED_OPERATORS_II}.

\paragraph*{} %- construction of generating operators
In the following we derive the generating operator
of a pair $(\vecV,\hilH)$ of Hilbert spaces in a straightforward way.
First of all,
each fixed $y \in \hilH$
gives rise to a linear functional on $\vecV$ according to
\[
f_y(v) := \scalar{\iota v}{y}_{\hilH}, \quad \quad (v \in \vecV).
\]
The mapping $f_y$ is also bounded,
{\ie} $f_y \in \vecV_a^*$.
Then,
by the \textsc{Riesz} representation theorem,
there is $v_y \in \vecV$ such that
\[
\scalar{v}{v_y} \= f_y(v) \qfa v \in \vecV.
\]
We therefore obtain a mapping $\linB\:\hilH \xto \vecV$,
by $\linB y := v_y$ for $y \in \hilH$,
which proves to be linear and continuous,
and the inner product in $\hilH$ is obtained from that in $\vecV$ by
\[
\scalar{\iota x}{y}_{\hilH} \= \scalar{x}{\linB y}_{\,\vecV}
\quad \quad (x,y \in \hilH).
\]
But the operator $\linB$ is also symmetric and strictly positive.
These properties imply the existence of the inverse operator
\[
\linB^{-1}\:\domain{\linB^{-1} } = \range{\linB} \subset \vecV \xto \hilH,
\]
which indeed is not necessarily bounded,
but uniquely defined by the pair $(\vecV,\hilH)$.

\subparagraph*{}
The operator $\linB^{-1}$ is symmetric and strictly positive,
too,
hence there exists the square root operator $\linB^{-1/2}$
having the property
\begin{equation}\label{GENERATING_OPERATOR}
\scalar{\linB^{-1/2} x}{\linB^{-1/2} y}_{\hilH} \= \scalar{x}{y}_{\,\vecV}
\quad \quad (x,y \in \hilH).
\end{equation}
It turns out that $\domain{\linB^{-1/2}} \eq \vecV$,
which can be shown by contradiction.
Especially for $y=x$ in \eqref{GENERATING_OPERATOR} we obtain the equality
\[
\snorm{\linB^{-1/2}x}_{\hilH} \= \norm{x}_{\,\vecV}
\quad \quad (x \in \vecV).
\]

%------------------------------------------------------------------------------%
\section{Sectorial Forms and Operators}                                        %
%------------------------------------------------------------------------------%
In this section we introduce
the notations of \textit{closed sectorial form} and
\textit{m-sectorial operator}
and show the relationships between them.

\subparagraph*{} %- sectorial forms
Let $\hilH$ be a complex Hilbert space
and $\sql\:\domain{\sql} \xtimes \domain{\sql} \xto \C$ be a sesquilinear form,
but defined on a linear subspace $\domain{\sql}$ of $\hilH$.
Like for linear operators,
the \textit{numerical range} of $\sql$ is defined by
\[
\Theta(\sql) \eqdef
\{\, \sql(x,x): x \in \domain{\sql}, \, \norm{x}=1 \,\}.
\]
We then denote the form $\sql$ to be \textit{sectorial},
if $\Theta(\sql)$ is contained in a sector $\Sigma_{\,\psi,\omega}$,
that is,
\begin{equation}\label{SECTOR_CONDITION}
\Theta(\sql) \, \subset \,
\Sigma_{\,\psi,\omega} \eqdef
\{\zeta \in \C: \zeta \neq \omega, arg(\zeta-\omega) < \psi \},
\end{equation}
where $\omega \in \R$ is the vertex
and $0 \leq \psi < \pi/2$ the semi-angle of the sector.
Notice that sectors are symmetric {\wrt} the real line.
Especially for $\psi \eq 0$ let $\Sigma_{\,0,\omega} := (\omega,\infty)$,
and for $\psi \eq \pi/2$ we have $\Sigma_{\,\pi/2,0} \eq \C_+$.
We also say that $\linA$ is $(\psi,\omega)$-sectorial,
if \eqref{SECTOR_CONDITION} holds.

\subparagraph*{} %- half-bounded forms
Now,
if $\sql$ is symmetric and $\sql \eq \sql^*$,
respectively,
then $\Theta(\sql) \subset \R$ and $\omega \leq \inf \Theta(\sql)$, 
so in this case we denote $\sql$ to be 
\textit{half-bounded} or \textit{bounded below}
and call $\omega$ the lower bound of $\sql$.
Moreover,
if $\sql$ is half-bounded and $(\psi,\omega)$-sectorial,
then the estimates
\[
\begin{cases}[1.2]
\quad \re \sql(x,x) \, \geq \, \omega \norm{x}^2 \\
\quad \im \sql(x,x) \, \leq \,
      \tan \psi \cdot (\re \sql(x,x) - \omega \norm{x}^2)\\
\end{cases}
\]
hold for all $x \in \domain{\sql}$.

\paragraph*{} %- closeness of sectorial forms
Let $\re \sql := (\sql + \sql^*)/2$ be the \textit{real part} of $\sql$.
Then for every $\alpha > 0$ the slightly modified form
$[\cdot,\cdot]_{\alpha} := \re \sql - \omega + \alpha$
constitutes an inner product on the linear space $\domain{\sql}$.
It turns out
that the norms $\abs{\cdot} _{\alpha}$
derived from these inner products $[\cdot,\cdot]_{\alpha}$
are equivalent to each other.
We especially take $\alpha=1$,
pay attention to the inner product
\[
[x,y]_1 \= \re \sql(x,y) - \omega \scalar{x}{y} + \scalar{x}{y},
\qfa x,y \in \domain{\sql}
\]
and distinguish between the two cases {\wrt} completeness of the norm
\[
\abs{x} _1 \eqdef {[x,x]_1}^{1/2}
\quad \quad (x \in \domain{\sql}).
\]
In both cases it is our goal to obtain a Hilbert space $\hilH_{\sql}$
with inner product $[x,y]_1$ and norm $\abs{x} _1$
for $x,y \in \hilH_{\sql}$,
together with a linear embedding $\iota\:\hilH_{\sql} \xto \hilH$,
which should be continuous if possible,
so $(\hilH_{\sql},\hilH)$ hopefully becomes into a pair of Hilbert spaces.

\begin{itemize}[leftmargin=5mm]
\item[1.]
The normed linear space $(\domain{\sql},\abs{\,\cdot\,} _1$)
is \textit{complete}.\\
We denote the form $\sql$ \textit{closed} in this case,
and $\domain{\sql}$ becomes into a Hilbert space in its own right,
hence we may define $\hilH_{\sql} := \domain{\sql}$ in this case. 
Notice that the norm $\abs{\,\cdot\,} _1$ on $\hilH_{\sql}$
is stronger than the norm $\norm{\,\cdot\,}$ in $\hilH$
restricted to $\hilH_{\sql}$,
so the natural inclusion mapping $\iota\: \hilH_{\sql} \xto \hilH$
proves to be continuous.

\item[2.]
The normed linear space $(\domain{\sql},\abs{\,\cdot\,} _1)$
is \textit{not} complete rsp. $\sql$ is not closed.\\
We define the Hilbert space $\hilH_{\sql}$ as the completion
of $\domain{\sql}$ {\wrt} the norm $\abs{\,\cdot\,} _1$.
Moreover,
Cauchy sequences in $\hilH_{\sql}$ are also Cauchy sequences in $\hilH$,
hence for every $x \in \hilH_{\sql}$
there is $x' \in \hilH$ such that $x' \eq \iota x$
for a linear map $\iota\:\hilH_{\sql} \xto \hilH$.

\subparagraph*{}
We want to extend the form $\sql$ from $\domain{\sql}$ to $\hilH_{\sql}$
as well,
and this will be done by defining
\[
\overline{\sql}(x,y) \eqdef \lim_{\,n \xto \infty} \sql(x_n,y_n)
\]
for $x,y \in \hilH_{\sql}$ and arbitrary sequences
$\{x_n\}_{n \in \N}, \{y_n\}_{n \in \N} \subset \domain{\sql}$
with $x_n \xto x$ and $y_n \xto y$.
It turns out that
the definition of $\overline{\sql}(x,y)$ is independent of
the sequences $\{x_n\}_{n \in \N}$ and $\{y_n\}_{n \in \N}$,
so this value is in fact well-defined for $x,y \in \hilH_{\sql}$.
\end{itemize}

\paragraph*{} %- closable sectorial forms
In summary,
we obtain in both cases a Hilbert space $\hilH_{\sql}$ with inner product
$[\cdot,\cdot]_1$ and norm $\abs{\,\cdot\,} _1$ together with a linear embedding
$\iota\: \hilH_{\sql} \xto \hilH$,
but which in the second case is not necessarily continuous.

\subparagraph*{} 
Regarding the second case,
the form $\sql$ is said to be \textit{closable},
if the mapping $\iota$ is continuous.
This definition especially means that every closed form is closable.
If $\sql$ is closable,
then the final result in both cases will be
that there exists a closed form,
namely $\sql$ itself in the first and $\overline{\sql}$ in the second case,
together with a continuous embedding of $\hilH_{\sql}$ into $\hilH$.

\paragraph*{} %- sectorialness
Let us continue with the concept of \textit{sectorialness}
of a closed linear operator in a complex Banach or Hilbert space.
This class of operators has been introduced and first studied
by \textsc{T.\,Kato}.
Let again $\Sigma_{\,\psi,\omega}$
be an open sector with vertex $\omega \in \R$ and semi-angle $0 < \psi \leq \pi$.
In his book \cite{VIII_Ka},
\textsc{Kato} denotes a linear operator $\linA$ in a Hilbert space
to be \textit{sectorial},
if the following conditions are fulfilled:

\begin{itemize}[labelsep=4mm,leftmargin=12mm]
\item[(S-1)]
$\linA$ is {quasi-accretive},
{\ie} $\scalar{(\linA + \omega) x}{x} \geq 0$
for some $\omega \in \R$ and $x \in \domain{\linA}$,

\item[(S-2)]
the numerical range $\Theta(\linA)$ is a subset of the sector
$\Sigma_{\,\psi,\omega}$.
\end{itemize}

Furthermore,
$\linA$ is called \textit{m-sectorial},
if it is at once sectorial and m-quasi-accretive,
that is,
$\linA + \omega$ is m-accretive for some $\omega \in \R$.

\paragraph*{} %- correspondence between closable sectorial forms and operators
We finally want to elaborate the correspondence
between closed sectorial forms and m-sectorial operators.
It should be clear
that if $\linA$ is a m-sectorial operator in $\hilH$
with domain of definition $\domain{\linA}$,
then the sesquilinear form
\[
\sql_{\linA}: \;
\begin{cases}[1.2]
\quad \domain{\sql_{\linA}} \eqdef \domain{\linA}, \\
\quad \sql_{\linA}(x,y) \eqdef \scalar{\linA \,x}{y}
\, \for x,y \in \domain{\sql_{\linA}} \\
\end{cases}
\]
is sectorial and closed.

\subparagraph*{}
Reversely,
let $\sql_0$ be a closable sectorial form in $\hilH$
with domain $\domain{\sql}$.
We consider the closed completion $\sql$ on the space $\hilH_{\sql}$.
A densely defined linear subspace $D \subset \domain{\sql}$ is denoted
a \textit{core} of $\sql$,
or equivalently,
the set $D$ is a core of $\sql$,
if $\sql$ is the closure of the restriction of $\sql$ to $D$.

\subparagraph*{}
\textsc{Kato}\textit{'s first representation theorem} claims that
a densely defined and closed sectorial form gives rise
to a m-sectorial linear operator:

\begin{theorem}\label{FIRST_REPRESENTATION_THEOREM}
Let $\sql$ be a densely defined and closed sectorial form in $\hilH$
with vertex $\omega$, semi-angle $\psi$ and domain of definition $\hilH_{\sql}$,
and let $\iota$ be the continuous embedding from $\hilH_{\sql}$ into $\hilH$.
Then there is a m-sectorial operator $\linA_{\sql}$ in $\hilH$,
whose domain of definition $\domain{\linA_{\sql}}$ is defined as follows:
Any element $x \in \hilH$ belongs to $\domain{\linA_{\sql}}$,
if $x \eq \iota x_0$ for some $x_0 \in \hilH_{\sql}$ and
$\sql(x_0,z) \eq \scalar{y}{\iota z}$ holds
for some $y \in \hilH$ and all $z \in \hilH_{\sql}$,
which suggests to define $\linA_{\sql} \, x := y$.

\paragraph*{}
The operator $\linA_{\sql}$ owns a couple of properties:

\begin{itemize}[leftmargin=5mm]
\item[(a)]
$\domain{\linA_{\sql}}$ is dense in $\hilH$
and the preimage $\iota^{-1} \domain{\linA_{\sql}}$
is dense in $\hilH_{\sql}$.

\item[(b)]
If $x \in \hilH_{\sql}$, $y \in \hilH$
and $\sql(x,y) \eq \scalar{\iota x}{y}$ for all $x$ in a core of $\sql$,
then $\iota x \in \domain{\linA_{\sql}}$ and $y \eq \linA (\iota x)$.

\item[(c)]
$\Theta(\linA_{\sql})$ is a dense subset of $\Theta(\sql)$.

\item[(d)]
The adjoint form $\sql^*$ is also closed and sectorial
with vertex $\omega$ and semi-angle $\psi$.

\item[(e)]
The adjoint ${\linA_{\sql}}^*$ of $\linA_{\sql}$
is the $m$-sectorial operator associated with the adjoint form $\sql^*$.
\end{itemize}
\end{theorem}

\subparagraph*{}
Thus m-sectorial linear operators are in one-to-one correspondence
with closed sectorial forms.

%------------------------------------------------------------------------------%
\section{Self-Adjoint Extensions of Half-Bounded Operators}                    %
%------------------------------------------------------------------------------%
We want to apply the results of the last section
to the following class of operators:
a \textit{symmetric} linear operator $\linA$ is called \textit{half-bounded},
if there is $\theta \in \R$ such that
\[
\scalar{\linA \,x}{x} \, \geq \, \theta \norm{x}^2
\qfa x \in \domain{\linA},
\]
or equivalently,
if $\Theta(\linA) \subset [\theta,\infty)$ for some $\theta \in \R$.
The operator $\linA$ is called \textit{positive} (\textit{strictly positive}),
if $\theta \eq 0$ ($\theta > 0$) can be chosen.
In case of $\theta > 0$,
the {Cauchy-Schwarz} inequality implies
that $\linA$ is also bounded below,
hence $\linA$ is surely injective.
In addition,
if $\linA$ is closed,
then the range $\range{\linA}$ is closed in $\hilH$,
which follows from the bounded-below theorem,
but $\range{\linA}$ may be a proper subspace of $\hilH$.

\paragraph*{}
The main result in this section is given by the fact
that each densely defined and half-bounded linear operator $\linA$
acting in a Hilbert space owns at least one \textit{self-adjoint extension}.

\subparagraph*{}
First notice that if $\linA$ is half-bounded with $\theta \in \R$,
then it is $(0,\theta)$-sectorial,
and the sesquilinear form $\sql(x,y) := \scalar{\linA x}{y}$
is sectorial and densely defined.
But $\sql$ is also closable,
so we may switch to the sesquilinear form $\overline{\sql}$ defined on the
Hilbert space $\hilH_{\sql}$,
called the \textit{energetic space} of the operator $\linA$.
Then by theorem \ref{FIRST_REPRESENTATION_THEOREM}
there is a m-sectorial and even self-adjoint linear operator $\linA_{\sql}$,
that forms an extension of the half-bounded operator $\linA$,
called the \textit{Friedrichs extension} of $\linA$.

\subparagraph*{}
The domain of definition $\domain{\linA_{\sql}}$ of the operator $\linA_{\sql}$
coincides with the linear subspace $\hilH_{\sql} \, \cap \, \domain{\linA^*}$.
Moreover,
the operator $\linA_{\sql}$ owns a useful property:
if one adds to the half-bounded operator $\linA$ the multiplication operator 
$x \in \hilH \mapsto \theta x$ with $\theta \in \R$,
then the {Friedrichs} extensions of $\linA$ and $\linA + \theta$
are related to each other by
$(\linA + \theta)_{\sql} \eq \linA_{\sql} + \theta$.

%------------------------------------------------------------------------------%
\section{Variational Evolution Equations}                                      %
%------------------------------------------------------------------------------%
Having the concept of Hilbert-Sobolev space available,
we want to review the variational approach to first and second order
evolution equations,
which has been developed by \textsc{J.-L.\,Lions}
in the midth of the last century.
It can be seen as the straightforward extension of
the \textsc{Lax-Milgram} theory from stationary equations to evolutional ones.

\subparagraph*{}
Let $\T > 0$ and denote $\lambda$ the Lebesgue-Borel measure on $[0,\T)$.
By an \textit{abstract variational problem of parabolic type}
in a real Hilbert space triple $(\vecV,\hilH,\vecV^*)$
we mean a $\lambda$-nearly everywhere in $(0,\T)$ valid
linear operator-differential equation
\begin{equation}\label{ABSTRACT_PARABOLIC_VARIATIONAL_PROBLEM}
\begin{cases}[1.2]
\quad u'(t) \, + \, \repA(t) \, u(t) \= f(t), \\
\quad u(0) = u_0, \\
\end{cases}
\end{equation}
where $u_0 \in \hilH$,
$f\:[0,\T) \xto \vecV^*$
and $\repA_{t \in [0,\T)}$ is a family
of linear operators $\repA\:\vecV \xto \vecV^*$.

\subparagraph*{} %- definition of weak solution I
Let $\sql(t)$ be the bilinear form associated with the operator $\repA(t)$.
Then we call $u \in H^1(0,\T;\vecV,\vecV^*)$
a \textit{weak solution} of \eqref{ABSTRACT_PARABOLIC_VARIATIONAL_PROBLEM},
if $u(0) \eq u_0$ and
\[
\scalar{u'(t)\,}{v(t)} \, + \, \sql(t,u(t),v(t)) \= f(t)(v)
\]
for all $v \in \vecV$ and $\lambda$-nearly everywhere $t \in [0,\T)$.

\subparagraph*{} %- Lions theorem for parabolic problems
Regarding the existence and uniqueness of a \textit{weak solution}
for the initial value problem \eqref{ABSTRACT_PARABOLIC_VARIATIONAL_PROBLEM},
the following theorem is of fundamental meaning:

\begin{theorem}\label{LIONS_THEOREM_I}
Let $\{\sql(t)\}_{t \in [0,\T)}$ fulfill the following conditions:

\begin{itemize}[leftmargin=9mm]
\item[(a)]
the mapping $t \in [0,\T) \mapsto \sql(t,u,v)$ is Lebesgue-measurable
for every fixed pair $(u,v) \in \vecV \xtimes \vecV$,

\item[(b)]
$\sql$ is \textit{uniformly} bounded in $[0,\T)$,
that is,
there is a constant $M \geq 0$ such that
\[
\abs{\sql(t,u,v)} \, \leq M \cdot \norm{u}_{\vecV} \norm{v}_{\vecV}
\qfa u,v \in \vecV, \; t \in [0,\T)\,,
\]
\item[(c)]
$\sql(t)$ is quasi-coercive for each $t \in [0,\T)$.
\end{itemize}

Then for each $u_0 \in \hilH$ and each $f \in L^2(0,\T;\vecV^*)$
there is exactly one solution $u \in H^1(0,\T;\vecV,\vecV^*)$
of \eqref{ABSTRACT_PARABOLIC_VARIATIONAL_PROBLEM},
which continuously depends on the data $u_0$ and $f$.
\end{theorem}

\subparagraph*{}
The proof of this theorem has been published in \cite{Li2}
as \textit{Th\'{e}or\`{e}me 1.1}, page 46.
Both terms $u'$ and $\repA(t) u$
in \eqref{ABSTRACT_PARABOLIC_VARIATIONAL_PROBLEM}
belong to the same space $L^2(0,\T;\vecV^*)$ as $f$,
which gives rise to say
that the problem has \textit{maximal regularity} in $\vecV^*$.

\paragraph*{} %- second order equations
The concept of weak solution makes also sense
for linear evolution equations of second order.

\subparagraph*{}
An \textit{abstract hyperbolic variational problem}
{\wrt} the Hilbert space triple $(\vecV,\hilH,\vecV^*)$ 
and a family of bilinear forms
$\sql(t)\:[0,\T] \xtimes \vecV \xtimes \vecV \xto \R$ is given by
\begin{equation}\label{ABSTRACT_HYPERBOLIC_VARIATIONAL_PROBLEM}
\begin{cases}[1.2]
\quad \scalar{u''(t)\,}{v(t)} + \sql(t,u(t),v(t)) = \scalar{f(t)}{v} 
\quad (v \in \vecV, \\
\quad u(0)  = u_0 \in \vecV, \\
\quad u'(0) = u_1 \in \hilH, \\
\end{cases}
\end{equation}
where $\scalar{\cdot}{\cdot}$ is the inner product in $\hilH$.

\paragraph*{}
The next theorem shows that
problem \ref{ABSTRACT_HYPERBOLIC_VARIATIONAL_PROBLEM}
is \textit{weakly solvable} under a couple of conditions
on the involved bilinear forms:

\begin{theorem}\label{LIONS_THEOREM_II}
Let $\{\sql(t)\}_{t \in [0,\T)}$ be a family
of bilinear forms on $\vecV$ such that

\begin{itemize}[leftmargin=9mm]
\item[(a)]
$t \mapsto \sql(t)$ is \textit{uniformly} bounded in $[0,\T)$,

\item[(b)]
$\sql(t)$ is $\hilH$-elliptic for each $t \in [0,\T)$,

\item[(c)]
$\sql(t)$ satisfies the symmetry condition
$\sql(t,u,v) \eq \sql(t,v,u)$
for all $u,v \in \vecV, \; t \in [0,\T)$,

\item[(d)]
the mapping $t \mapsto \sql(t,u,v)$ is continuously differentiable in $[0,\T)$
for each fixed $(u,v) \in \vecV \xtimes \vecV$,

\item[(e)]
there is $M_1 \geq 0$ such that
\[
\abs{ \frac{d}{dt} \sql(t,u,v)} \, \leq \, 
M_1 \cdot \norm{u}_{\vecV} \norm{v}_{\vecV}
\qfa \, t \in [0,\T).
\]
\end{itemize}
Then the initial value problem \eqref{ABSTRACT_HYPERBOLIC_VARIATIONAL_PROBLEM}
has a unique weak solution $u \in H^1(0,\T;\vecV,\vecV^*)$
for all $u_0 \in \vecV$,
        $u_1 \in \hilH$ and
        $f \in L^2(0,\T;\hilH)$.
\end{theorem}

\subparagraph*{}
The symmetry condition (c) in \ref{LIONS_THEOREM_II} is essential
to prove this theorem.
Moreover,
the solution map
\[
R\:L^2(0,\T;\hilH) \xtimes \vecV \xtimes \hilH \xto L^2(0,\T;\hilH)
\]
given by $(f,u_0,u_1) \mapsto u$ is linear in each argument and continuous,
thus problem \eqref{ABSTRACT_HYPERBOLIC_VARIATIONAL_PROBLEM} proves
to be well-posed.

%------------------------------------------------------------------------------%
\section{Applications}                                                         %
%------------------------------------------------------------------------------%
In this final section we want to illustrate the usage of abstract sesquilinear
forms in Hilbert space for the analysis of linear equations of the form
$\A u \eq f$,
where $\A$ is a linear partial differential operator and $u,f$ are real
or complex-valued mappings defined on an open set $\G$
within an Euclidean space $\Rd$.

\paragraph*{} %- differential expressions in divergence form
A differential expression $\A$ of order $2m$
is said to be in \textit{divergence form},
if there are functions $a_{\alpha \beta} \in C^{\abs{\alpha} }(\G)$
for $\abs{\alpha} , \abs{\beta} \leq m$
such that $\A$ can be written as
\begin{equation}\label{VIII_LPDO_DIV}
\A \eqdef
\sum_{\abs{\alpha} \leq m, \abs{\beta} \leq m} (-1)^{\abs{\alpha} } \,
\PA(\,a_{\alpha \beta}(x) \, D^{\beta}).
\end{equation}

\subparagraph*{}
In opposite to this definition we call an arbitary differential expression
of order $2m$ given by
\begin{equation}\label{VIII_LPDO}
\A'(x,\,\cdot\,) \eqdef
\sum_{\abs{\alpha} \leq 2m} c_{\alpha}(x) \, \PA
\end{equation}
to be \textit{in non-divergence form}.

\subparagraph*{}
It is a matter of fact
that every linear differential expression $\A$ in divergence form
can be transformed into an expression $\A'$ in \textit{non-divergence form}.
Reversely,
a non-divergence form expression $\A'$ given by \eqref{VIII_LPDO}
can be transformed into an expression $\A$ in divergence form \eqref{VIII_LPDO_DIV},
if $c_{\alpha} \in C^{2m - \abs{\alpha} }(\G)$ for $\abs{\alpha} \leq 2m$,
that is,
if the coefficients are sufficiently many differentiable in $\G$.
By comparing the expressions \eqref{VIII_LPDO} and \eqref{VIII_LPDO_DIV},
it turns out that
$\A'$ is strongly elliptic in $\G$ \iff $\A$ is so.

\subparagraph*{}
Next we elaborate the purpose of a differential operator
being in divergence form.
Let $\TF$ be the space of \textit{test functions}
and furnished with the inner product
\[
\scalar{\phi}{\psi} \eqdef
\int_{\,\G} \phi \, \overline{\psi} \,dx \qfa \phi,\psi \in \TF.
\]
Now,
applying $\A$ given in divergence form to $\phi \in \TF$,
then multiplying $\A \, \phi$ with $\overline{\psi} \in \TF$ and
finally performing $m$-times partial integration leads to
\begin{equation}\label{DIRICHLET_FORM}
\DF_{\A}(\phi,\psi) \eqdef \scalar{\A \, \phi}{\psi} \=
\sum_{\abs{\alpha} \leq m, \abs{\beta} \leq m} \,
\int_{\G}
a_{\alpha \beta} \; \PA \phi \; D^{\beta} \overline{\psi}
\, dx,
\end{equation}
since all the boundary terms vanish.
The mapping
\[
\DF_{\A} := \TF \xtimes \TF \xto \C
\]
is a sesquilinear form,
called the \textit{Dirichlet form} of $\A$ in $\G$.
If the coefficent functions $a_{\alpha \beta}$ of $\A$
are measurable and essentially bounded,
then the mapping $\DF_{\A}$ is continuous
and thus can continuously be extended to the product space
$\HN{m}{\G} \xtimes \HN{m}{\G}$,
where $\HN{m}{\G}$ is the Sobolev space with null boundary conditions,
that is,
the completion of $\TF$ {\wrt} the Sobolev norm $\abs{\,\cdot\,} _m$.

\paragraph*{}
We are mainly interested in a $2m$-order linear differential operator $\A$
given in divergence form,
that together with the underlying domain $\G$ fulfills the following conditions:

\begin{itemize}[leftmargin=14mm]
\item[(I)\;\;]
$\G$ is bounded in $\Rd$,

\item[(I')\;\,]
$\G$ satisfies the $m$-th \textit{extension property},
that is,
the linear space $R^0(\G)$ consisting of those functions $\phi \in W^{m,p}(\G)$,
which are restrictions of functions $\phi_{ext} \in C_c^{\infty}(\Rd)$ to $\G$,
is dense in $W^{m,p}(\G)$.

\item[(II)\;]
$\A$ is uniformly strongly elliptic in $\G$,
that is,
there is $\theta > 0$ %(independent of $x \in \clG$)
such that
$\re \, (-1)^m \, p(x,\xi) \geq \theta \cdot \abs{\xi} ^{2m}$
for all $\xi \neq 0$ and $x \in \clG$.

\item[(III)\,]
the coefficient functions $a_{\alpha \beta}$ are essentially bounded in $\G$,
that is,
$a_{\alpha \beta} \in L^{\infty}(\G)$,

\item[(III')]
the coefficient functions $a_{\alpha \beta}$ of highest order
$\abs{\alpha} + \abs{\beta} \eq 2m$ are uniformly continuous in $\clG$.
\end{itemize}

\subparagraph*{}
\textsc{{\GARDING}}\textit{'s inequality} is the main assertion in this section
and can be seen as a bridge
from the theory of linear-elliptic differential operators
to abstract functional analysis:

\begin{theorem}\label{GARDING_INEQUALITY_THEOREM}
Let the domain $\G$ fulfill the conditions (I),\,(I') 
and let the operator $\A$ of order $2m$
fulfill the conditions (II),\,\,(III),\,(III').
Then there is $\lambda_0 \in \R$ and $a_0 > 0$ such that
for all $u \in \HN{m}{\G}$
\begin{equation}\label{GARDING_INEQUALITY}
\DF_{\A}(u,u) \, + \, \lambda_0 \norm{u}^2
\, = \,
\scalar{\,(\A+\lambda_0)\,u}{u}
\, \geq \,
a_0 \, {\abs{u} _m}^2.
\end{equation}
\end{theorem}

\subparagraph*{}
In other words,
\textsc{{\GARDING}}'s inequality quotes that
under the mentioned conditions the Dirichlet form associated with
the operator $\A + \lambda_0$ is coercive,
hence the \textsc{Lax-Milgram} theorem may be applied.

\subparagraph*{} %- Garding's inequality for second order operators
% We do not want to prove theorem \ref{GARDING_INEQUALITY_THEOREM}
% and refer to \cite{VIII_RR},\,p.\,292,
% but instead will prove the version for differential operators
% of \textit{second order},
The most important class of linear partial differential operators
in divergence form {\wrt} applications is that of \textit{second order},
that is,
our operator $\A$ is of the form
\begin{equation}\label{GENERAL_SECOND_ORDER_DIFFERENTIAL_OPERATOR}
\A(x,\cdot) \= 
-\sum_{i=1}^d \,
\partial_i \, (\sum_{j=1}^d a_{ij}(x)) 
\, + \,
\sum_{j=1}^d \, (b_j(x) \, \partial_j)
\, + \,
c(x),
\end{equation}
which is nothing but \eqref{VIII_LPDO_DIV} for $m=1$.

\subparagraph*{} %- Dirichlet form
The Dirichlet form $\DF_{\A}$ belonging to $\A$ 
given by \eqref{GENERAL_SECOND_ORDER_DIFFERENTIAL_OPERATOR}
is the sesquilinear form defined by
\begin{equation}\label{GENERAL_SECOND_ORDER_DIRICHLET_FORM}
\DF_{\A}(u,v) \= \int_{\G} \,
(\sum_{i,j=1}^d a_{ij} \, \partial_i u \, \partial_j \overline{v} \, + \,
 \sum_{j=1}^d b_j \, \partial_j u \, \overline{v} \, + \,
 c \,u \, \overline{v}
)
\, dx.
\end{equation}
In this case,
\textsc{\GARDING}'s inequality remains even true
if condition (I)' and (III') are dropped:

\begin{theorem}\label{GARDING_INEQUALITY_THEOREM_II}
Let $\G \subset \Rd$ fulfill (I) and
the differential operator $\A$ of second order and given in divergence form
fulfill (II) and (III).
Then there is $\lambda_0 \in \R$ and $a_0 > 0$ such that
\[
\DF_{\A}(u,u) \, + \, \lambda_0 \norm{u}^2
\=
\scalar{\,(\A+\lambda_0)\,u}{u}
\, \geq \,
a_0 \, {\abs{u} _1}^2
\]
for all $u \in \HN{1}{\G}$,
where $a_0$ is the uniform lower bound of the eigenvalues of all the matrices
$A(x)$ with $x \in \G$.
\end{theorem}

\subparagraph*{} %- proof of Garding's inequality for second order operators
To prove \textsc{\GARDING}'s inequality in the special case of second order operators,
first recall \textsc{Young}\textit{'s inequality}
$ab \leq \epsilon a^2 + 4b^2/ \epsilon$
for $a,b \in \R$ and any $\epsilon >0$.
By the strict ellipticity condition for $\A$ and with
$b_0 := \max_{1 \leq j \leq d} \norm{b_j}_{\infty}$
and $c_0 := \norm{c}_{\infty}$
we have
\[
\DF_{\A}(u,u) \= \int_{\,\G} \,
(
 \sum_{i,j=1}^d a_{ij} \, \partial_i u \, \partial_j \overline{u} \, + \,
 \sum_{j=1}^d   b_j \, \partial_j u \, \overline{u} \, + \,
 c \, u \, \overline{u}
) \, dx
\]
\[
\geq \, 
\theta \int_{\,\G} \abs{\nabla u} ^2 \, dx
\, - \,
b_0 \int_{\,\G} \abs{\nabla u} \, \abs{u} \, dx
\, - \, c_0 \norm{u}^2
\]
\[
\geq \, 
\theta \int_{\,\G} \abs{\nabla u} ^2 \, dx
\, - \,
b_0 \epsilon           \int_{\,\G} \abs{\nabla u} ^2 \, dx
\, - \,
\frac{b_0}{4 \epsilon} \int_{\,\G} \abs{u}        ^2 \, dx
\, - \, 
c_0 \norm{u}^2
\]
for some $\epsilon > 0$.
Let us chose $\epsilon$ small enough to ensure 
$
b_0 \, \epsilon \leq \frac{\theta}{2},
$
then
\[
\DF_{\A}(u,u)
\, \geq \, 
\frac{\theta}{2} \int_{\,\G} \abs{ \nabla u} ^2 \, dx
 \, - \, \,
(\frac{b_0}{4 \epsilon} + c_0) \norm{u}^2
\=
\]
\[
\frac{\theta}{2} (\int_{\,\G} (\abs{u} ^2 + \abs{ \nabla u} ^2) \, dx
 \, - \, \,
(\frac{b_0}{4 \epsilon} + c_0 + \frac{\theta}{2}) \norm{u}^2,
\]
hence with $a_0 := \theta/2$  and
$\lambda_0 := b_0 / 4 \epsilon + c_0 + \theta/2$
we obtain \eqref{GARDING_INEQUALITY}.

\paragraph*{} %- coerciveness of second-order differential operators
Recall again
that a sesquilinear form $\sql$ acting in a complex Hilbert space $\vecV$
is coercive,
if $\re \sql(x,x) \geq \theta \norm{x}^2$ holds for all $x \in \vecV$.
In the following we again consider the Dirichlet form $B_{\A}$
of a second-order differential operator $\A$
given in divergence form and 
by \eqref{GENERAL_SECOND_ORDER_DIFFERENTIAL_OPERATOR}.
A short computation shows
that $\DF_{\A}$ is coercive,
if the function $c$ is large enough in comparison
with the sum of partial derivatives of the functions $b_k$.
More precisely,
$B_{\A}$ proves to be coercive if
\begin{equation}\label{COERCIVITY_CONDITION}
c(x) \, \geq \, \frac{1}{2} \,
\sum_{k=1}^d \, \frac{\partial b_k}{\partial x_k}(x)
\quad \quad (x \in \G).
\end{equation}

\subparagraph*{}
On the other hand,
if the matrix $A(x) \eq (a_{ij}(x))_{i,j=1}^d$ consistent of the second-order
coefficients of $\A$ fails to be positive definite for all $x \in \G'$,
where $\G' \subset \G$ has positive measure,
then $\DF_{\A}$ cannot be coercive on $\HN{1}{\G}$.
Otherwise,
if $A(x)$ is positive definite in all $x \in \G$ and 
condition \eqref{COERCIVITY_CONDITION} does not hold.
then the functions $b_j$ disturb the coercivity of the form $\DF_{\A}$.
But after all,
the form $\DF_{\A}$ is at least quasi-coercive,
which is nothing but the assertion of \textsc{\GARDING}'s inequality.

\paragraph*{} %- Gardings inequality for the Laplacian
Let us especially consider the Dirichlet form $B_{-\Delta}$
associated with the negative Laplacian $-\Delta$,
which adopts the simple form
\begin{equation}\label{LAPLACE_DIRICHLET_FORM}
\DF_{-\Delta}(u,v) \= \int_{\,\G} \nabla u \cdot \nabla \overline{v} \, dx
\end{equation}
for $u,v \in \HN{1}{\G}$.
\textsc{{\GARDING}}'s inequality for this form reads as 
\[
\DF_{-\Delta}(u,u) \, + \, \norm{u}^2 \=
\int_{\,\G} (\abs{\nabla u} ^2 + \abs{u} ^2) \, dx \= {\abs{u} _1}^2,
\]
so $a_0 \eq \lambda_0 \eq 1$ can be chosen in this case.

\subparagraph*{} %- Poincare's inequality
The assertion that $\DF_{-\Delta}$ is coercive in bounded domains,
is of course a consequence of \textsc{{\GARDING}}'s inequality,
but may also proved directly by \textsc{{\POINCARE}}\textit{'s inequality}
even for non-bounded domains.
Let the mapping $\abs{\,\cdot\,} _e \:\HN{1}{\G} \xto [0,\infty)$ be defined by
\[
\abs{u} _e \eqdef
\left(
\int_{\G} \nabla u \, \nabla \overline{u} \; dx
\right) ^{1/2},
\]
which fulfills all the properties of a norm.
It is denoted the \textit{energetic norm} in $\HN{1}{\G}$.
Suppose that $\G$ has at least a bound in one direction,
say $\abs{x_1} \leq \delta$
for all $x=(x_1,\ldots,x_d) \in \G$ and some $\delta > 0$.
Then there is a constant $C(\G) > 0$
such that
\begin{equation}\label{POINCARE_INEQUALITY}
\abs{u} _1 \, \leq \, C(\G) \cdot {\abs{u} _e}
\end{equation}
for all $u \in \HN{1}{\G}$,
where the lower bound $C(\G) > 0$ only depends on the shape of $\G$.
This inequality claims
that in $\HN{1}{\G}$
the energetic norm $\abs{\,\cdot\,} _e$
and the Sobolev norm $\abs{\,\cdot\,} _1$
are equivalent to each other.
\begin{comment}
For the proof of \eqref{POINCARE_INEQUALITY},,
we perform integration by parts and obtain
\[
\norm{u}_2^2
\, = \,
\int_{\G} 1 \cdot \abs{u} ^2 \, dx
\, = \,
- \int_{\G} x_1 \, \frac{\partial}{\partial x_1} \abs{u} ^2 \, dx
\, \leq \,
2 \delta \norm{u}_2 \norm{\partial_1 u}.
\]
\end{comment}

\subparagraph*{}
\textsc{{\POINCARE}}'s \textit{inequality} also implies that
the form $\DF_{-\Delta}$ as defined in \eqref{LAPLACE_DIRICHLET_FORM}
is coercive.
The linear operator $\linA_{\sql}$ associated with $\sql \eq \DF_{-\Delta}$
has domain of definition
\[
\domain{\linA_{\sql}} \= \{\,u \in \HN{1}{\G}: \linA_{\sql} u \in \Lh\,\}
\]
with values $\linA_{\sql} u := -\Delta u$ in the sense of weak derivatives.
This operator is in fact the realization of the negative Laplacian
in Hilbert space $\Lh$ and
under homogeneous Dirichlet conditions.
By theorem \ref{QUASI-COERCIVE_FORMS_AND_ASSOCIATED_OPERATORS},
$\linA_{\sql}$ is a densely defined and closed operator.

\paragraph*{} %- the variational Dirichlet problem
Let us now turn back to higher order differential operators and
formulate the existence and uniqueness assertion
for stationary equations of the form $\A\,u=f$,
where

\begin{itemize}[leftmargin=15mm]
\item[(DP-1)]
the differential operator is of order $2m$,
strongly elliptic in a \textit{bounded} domain $\G \subset \Rd$
and given in divergence form,
hence
\[
\A \eqdef
\sum_{\abs{\alpha} \leq m, \abs{\beta} \leq m} (-1)^{\abs{\alpha} } \,
\PA(\,a_{\alpha \beta}(x) \, D^{\beta}),
\]

\item[(DP-2)]
the coefficients fulfill $a_{\alpha \beta} \in L^{\infty}(\G)$ and
the top-order coefficients are uniformly continuous in $\clG$.
\end{itemize}

\paragraph*{} %- homogeneous Dirichlet problem
Our aim is to solve the \textit{homogeneous Dirichlet problem}
\[
\begin{cases}[1.2]
\quad (\A\,u)(x)\= f(x)& (x \in \G),\\
\quad u(x) =
\partial u/ \partial \nu (x) =
\ldots = \partial u^{m-1} / \partial \nu^{m-1}(x) = 0 & (x \in \bndG).\\
\end{cases}
\]

\subparagraph*{}
Let $\vecV \eq \HN{m}{\G}$ and $\hilH \eq \Lh$,
so the pair $(\vecV,\hilH)$ forms a Hilbert triple.
By \textsc{{\GARDING}}'s inequality \eqref{GARDING_INEQUALITY},
the Dirichlet form $\DF_{\A}(u,v) := \scalar{\A u}{v}$ for $u,v \in \vecV$
of $\A$ is quasi-coercive with some $\lambda_0 \in \R$.
Then the \textsc{Lax-Milgram} theorem implies
the \textit{existence and uniqueness theorem}
for the homogeneous Dirichlet problem:

\begin{theorem}\label{DIRICHLET_EXISTENCE_THEOREM}
For $\lambda \geq \lambda_0$, the variational equation
\begin{equation}\label{VARIATIONAL_EQUATION_II}
\DF_{\A}(u,v) \, + \, \lambda \scalar{u}{v}
\, = \, f(v) \quad (v \in \HN{m}{\G})
\end{equation}
owns for each $f \in H^{-m}(\G)$
a unique solution $u \in \HN{m}{\G}$.
\end{theorem}

\subparagraph*{} %- weak solution of the Dirichlet problem
The function $u$ is called the \textit{weak solution}
of the homogeneous Dirichlet problem.

\paragraph*{} %- notes on elliptic regularity
Having theorem \ref{DIRICHLET_EXISTENCE_THEOREM} available,
it is the task of \textit{elliptic regularity theory} to show,
that if $\G$ is bounded,
$\bndG$ is sufficiently smooth
and $f \in H^k(\G)$ for some $k \in \N_0$,
then the solution $u$ of \eqref{VARIATIONAL_EQUATION_II}
actually belongs to the space
$H^{k+2m}(\G) \, \cap \, \HN{m}{\G}$.

\subparagraph*{}
Moreover,
let us assume
that the existence of a weak solution
for equation $\A\,u=f$
under homogeneous Dirichlet boundary conditions is ensured.
Then the question arises,
under which conditions the weak solution $u \in \HN{m}{\G}$
is classical,
that is,
$u$ is contained in the space $C^{2m}(\G) \, \cap \, C^m(\clG)$.

\paragraph*{}
In the following
we present only few results on regularity theory
and restrict ourselves to the second order differential operator
\[
\A_2(x,u) \=
- \sum_{i,j=1}^d \partial_i (a_{ij}(x) \, \partial_j u(x))
+
\sum_{j=1}^d b_j(x) \, \partial_j u(x)
+
c(x) \,u(x).
\]
Basically,
we distinguish between
\textit{inner regularity} and \textit{regularity up to the boundary}.

\paragraph*{} %- inner regularity
Starting with \textit{inner regularity},
let $\A_2$ be strongly elliptic and of second order,
whose coefficients fulfill
$a_{ij} \in C^1(\G)$, $b_i \in L^{\infty}(\G)$ and $c \in L^{\infty}(\G)$.
If $f \in \Lh$ and $u \in H^1(\G)$ is a weak solution of $\A_2 \, u \eq f$,
then the following propositions are valid:

\begin{itemize}
\item[1.]
$u \in H_{loc}^2(\G)$,
that is,
for each open subset $\G' \Subset \G$\footnote{
\,$\G' \Subset \G$ means
that $\G'$ is \textit{compactly embedded} in $\G$,
{\ie}
there is a compact set $\compK$ such that $\G' \subset \compK \subset \G$.}
the function $u \mid_{\G'}$ belongs to the space $H^m(\G')$,

\item[2.]
for each open subset $\G' \Subset \G$ the estimate
\[
\norm{u}_{H^2(\G')} \, \leq \,
C \cdot [\,\norm{f}_{\Lh} + \norm{u}_{\Lh}\,]
\]
holds,
where the constant $C$ only depends on $\G'$, $\G$
and the coefficients of the operator $\A_2$.
\end{itemize}

Notice that
there are no assumptions on smoothness of the boundary $\bndG$ required.

\paragraph*{} %- higher inner regularity
A similar result is obtained,
called \textit{higher inner regularity},
if one requires higher smoothness properties for the function $f$ and
the coefficients of $\A_2$.
Let $a_{ij}, b_i, c \in C^{m+1}(\G)$, $f \in H^m(\G)$ for some $m \in \N$.
If $u \in H^1(\G)$ is a weak solution of the equation $\A_2 u =f$,
then we have $u \in H_{loc}^{m+2}(\G)$ and
\[
\norm{u}_{H^{m+2}(\G')} \, \leq \,
C \cdot [\, \norm{f}_{H^m(\G)} + \norm{u}_{\Lh} \,] \,,
\]
where $C$ again only depends on the sets $\G'$, $\G$
and the coefficients of $\A_2$.
Finally,
if $u \in H^1(\G)$ is the weak solution of $\A_2 \, u \eq f$
and $a_{ij}, b_i, c, f \in \SF$,
then we even have $u \in \SF$.

\paragraph*{} %- regularity up to the boundary
We now turn to \textit{regularity up to the boundary}.
This notation corresponds to assertions about the membership
of the weak solution to one of the Sobolev spaces $H^m(\G)$.
However,
to obtain this kind of regularity,
one has to postulate sharper requirements
on the coefficients of the differential operator
and certain smoothness properties of the boundary of $\G$.

\subparagraph*{}
This kind of regularity means
that if $k \in \N_0$,
$a_{ij}, b_j, c \in C^k(\clG)$,
$f \in H^k(\G)$,
the boundary $\bndG$ is of class $C^{k+2}$ and
$u \in \HN{1}{\G}$ is a weak solution of the homogeneous Dirichlet problem,
then $u \in H^{k+2}(\G)$ and the following Schauder estimate is valid:
\[
\norm{u}_{H^{k+2}(\G)} \, \leq \,
C \cdot [\, \norm{f}_{H^k(\G)} + \norm{u}_{\Lh} \,].
\]

\subparagraph*{}
There are statements about higher regularity of $u$ as well:
let $a_{ij}, b_i, c \in C^{\infty}(\clG)$,
$f \in \SF \, \cap \, \Lh$ and $\bndG$ be of class $C^{\infty}$,
then each solution $u \in H^1(\G)$
of the homogeneous Dirichlet problem is indefinitely differentiable
(\textsc{Weyl}\textit{'s lemma}).

\paragraph{} %- the Neumann Problem for the negative Laplacian
Besides the Dirichlet boundary value problem,
the \textit{Neumann problem} is the second important boundary value problem
for elliptic differential operators of second order
and especially for the negative Laplacian,
which in the following we shall sketch in the setting of Hilbert-Sobolev spaces.

\subparagraph*{} %- Neumann problem for the negative Laplacian
Let $\G \subset \Rd$ be a bounded domain with $C^1$-boundary,
so the outer normal vector $n(x)$ is defined in every point $x \in \bndG$.
Then the \textit{homogeneous} Neumann problem reads as follows:
find a function $u \in C^2(\G) \cap C^1(\overline{\G})$ such that
\begin{equation}\label{NEUMANN_PROBLEM}
\begin{cases}[1.2]
\quad (-\Delta u)x \= f(x), & x \in \G, \\
\quad u_{\nu}(x) := (\nabla u)(x) \cdot n(x) \= 0, & x \in \bndG \\
\end{cases}
\end{equation}
for given $f \in C(\G)$.
Notice that with each solution $u$ of this problem
the function $u_c := u+c$ ($c \in \R$) is a solution,
too.

\subparagraph*{} %- divergence theorem
The \textit{divergence theorem} quotes that
\[
\int_{\G} \Delta u \, dx \= \int_{\bndG} \nabla u \cdot n \, d\sigma
\]
holds for every function $u \in C^2(\G) \cap C(\bndG)$,
which implies that
a necessary condition for the existence of a solution
of \eqref{NEUMANN_PROBLEM} is the validness of the equation
\begin{equation}\label{NEUMANN_COMPATIBILITY_CONDITION}
\int_{\G} f dx \= -\int_{\G} \Delta u \, dx.
%\= \int_{\bndG} g \, d\sigma.
\end{equation}
Conversely,
if this compatiblity condition does not hold,
there cannot exist a solution of problem \eqref{NEUMANN_PROBLEM}.

\paragraph*{} %- principles of the variational approach
The variational approach to the Neumann problem consists of 
multiplying each side of $-\Delta u \eq f$ with functions $v \in H^1(\G)$ and
performing partial integration on the left hand side,
which leads to the variational equation
\[
\int_{\G} \nabla u \cdot \nabla v \, dx \=
\int_{\G} f \, v \, dx.
% \, - \, \int_{\bndG} g \, v \, d\sigma.
\]

\subparagraph*{}
Let $\gamma_0$ be the \textit{trace operator} on the Sobolev space $H^1(\G)$
into the space $L^2(\bndG)$. 
Then $u \in T_g(\G) := \{\, \phi \in H^1(\G): \gamma_0 \phi \eq 0 \,\}$,
the \textit{trial space} for the Neumann problem,
is called a \textit{weak solution} of \eqref{NEUMANN_PROBLEM},
if for all $v \in H^1(\G)$
\[
\int_{\G} \nabla u \cdot \nabla v \, dx \=
\int_{\G} f \, v \, dx.
%\, - \, \int_{\bndG} g \, \gamma_0 v \, d\sigma.
\]

\paragraph*{}
If a linear differential operator $\A$ is given in divergence form
and its Dirichlet form is coercive,
then the variational approach to the heat equation associated with $\A$
may be applied.
Consider for $f \in L^2([0,\T), H^{-1}(\G))$ and $u_0 \in \Lh$
the \textit{parabolic} initial-value problem
\[
\begin{cases}[1.2]
\quad \scalar{u_t}{v} \, + \, \DF_{\A}[u,v] \= \scalar{f}{v} 
\qfa v \in \HN{1}{\G},\\
\quad u(t,\cdot) \in \HN{1}{\G}, \\
\quad u(0) = u_0.\\
\end{cases}
\]
This problem owns a uniquely defined weak solution
in the space $H^1(0,\T; \HN{1}{\G}, H^{-1}(\G))$,
which fulfills the variational equation $\lambda$-nearly everywhere in $[0,\T)$.
Notice that this assertion is especially true if $\A = -\Delta$.

\subparagraph*{} %- variational approach to the hyperbolic problem
Again,
if $\A$ is given in divergence form,
the \textit{hyperbolic} initial-value problem in variational form
is given by
\[
\begin{cases}[1.2]
\quad \scalar{u_{tt}}{v} \, + \, \DF_{\A}[u,v] \= \scalar{f}{v} 
\qfa v \in \HN{1}{\G}, \\
\quad u(t,\cdot) \in \HN{1}{\G},\\
\quad u(0)=u_0, \\
\quad u_t(0)=u_1, \\
\end{cases}
\]
where the coefficients $a_{ij}$ of the operator $\A$ must fulfill
the symmetry conditions
\[
a_{ij}(x) \=
\overline{a_{ji}(x)} \quad 1 \leq i,j \leq d, \;\; x \in \G.
\]

\subparagraph*{}
If $f$ belongs to $L^2([0,\T), \Lh)$,
$u_0 \in \HN{1}{\G}$ and $u_1 \in \Lh$,
then this problem has exactly one weak solution $u \in L^2(0,\T; \HN{1}{\G})$
such that $u(0) = u_0$, $u_t(0)=u_1$
and $u(t) \in \HN{1}{\G}$ for $t \in [0,\T)$ $\lambda$-nearly everywhere.

\subparagraph*{}
Since the Dirichlet form of the expression $-\Delta$ is symmetric,
the variational initial-value problem for the classical wave equation
$u'' - \Delta u = f$
owns a uniquely defined weak solution
$u \in H^1(0,\T; \HN{1}{\G}, H^{-1}(\G))$,
if the functions $f, u_0, u_1$ fulfills the mentioned prerequisites.

%------------------------------------------------------------------------------%
%                                 BIBLIOGRAPHY                                 %
%------------------------------------------------------------------------------%
%\newpage
%\thispagestyle{empty}
\vspace{10mm}

\titleformat{\section}
{\normalfont \bfseries \filcenter}
{\thesection.\;}{0mm}{}

\begin{thebibliography}{99}
\begin{small}
\bibitem{VIII_Ka} T.\,Kato:
\textit{Perturbation Theory for Linear Operators}, Springer Verlag,
Grundlehren der mathematischen Wissenschaften, Band 132, 1995.

\bibitem{Li2} J.-L.\,Lions:
\textit{\'{E}quations Diff\'{e}rentielles Op\'{e}rationelles et Probl\`{e}mes aux Limites}.
Springer,
Grund\-lehren der mathematischen Wissenschaften, Band 111, 1961.

\bibitem{VIII_RR}  M.\,Renardy, R.\,Rogers:
\textit{An Introduction to Partial Differential Equations}.
Texts in Applied Mathematics 13, Springer, 2004.

\bibitem{VIII_Sh} R.\,Showalter:
\textit{Hilbert Space Methods in Partial Differential Equations},
Dover Publications, 2010.
\end{small}
\end{thebibliography}

%------------------------------------------------------------------------------%
%                               END OF DOCUMENT                                %
%------------------------------------------------------------------------------%
\end{document}
